\chapter{Conclusion and Outlook}
\label{chap:conclusion}

\section{Summary of Findings}
\label{sec:conc_summary}

This thesis investigated how architectural depth in image-domain cascaded encoder--decoder networks affects \gls{ct} restoration quality under compound corruption. Two contributions support this study. First, \Cref{sec:meth_overview} introduced a physics-informed multi-artifact sinogram corruption protocol that combines low-dose Poisson quantum noise with three independently gated physical artifact models: sinogram-domain motion, per-column ring offsets, and forward-projected metal masks with photon-starvation clipping. The protocol is implemented as a reproducible CPU-side dataset transform compatible with the \gls{lodopab} train/validation/test split. Second, \Cref{chap:experiments} presented a systematic comparison of independent-stage, end-to-end, and residual image-domain cascades against classical and learned single-stage baselines, together with per-stage loss and gradient analyzes used to interpret observed behavior.

The main empirical finding is that, under the compound corruption studied here, a single well-designed image-domain encoder--decoder already operates close to an effective performance ceiling that two-stage cascades do not surpass by depth alone. Concretely:

\begin{itemize}
    \item The large single-stage U-Net achieves \SI{36.11}{\decibel} PSNR and \num{0.8772} SSIM on the LoDoPaB-CT test split under the full compound-corruption distribution.
    \item At matched small per-stage capacity, the gain over the small single-stage U-Net increases monotonically across the cascade variants, from naive independent training to residual prediction to residual prediction with dual-filter input ($+0.15 \to +0.54 \to +0.72$\,dB).
    \item The strongest small-capacity cascade, namely the residual variant with dual-filter conditioning, matches the PSNR of the large single-stage U-Net while remaining slightly below it in SSIM.
    \item The end-to-end cascade exhibits progressive attenuation of the Stage~2 learning signal as Stage~1 approaches convergence, which helps explain why end-to-end coupling does not outperform the independently trained residual formulation.
    \item The asymmetric-capacity probe indicates that the second stage is constrained more by the residual information available to it than by its parameter count alone.
\end{itemize}

This is best interpreted as an architectural ceiling result rather than a negative one. The central conclusion is not simply that deeper image-domain cascades help little, but that additional depth becomes useful only when the later stage receives information that is not already compressed into the Stage~1 output.

\section{Author's Contributions}
\label{sec:conc_originality}

Following the \gls{lme} thesis guidelines (§1.7), the original contributions of this thesis are stated explicitly here. The author designed and implemented the physics-informed multi-artifact corruption protocol of \Cref{sec:meth_overview}, including the metal forward-projection library and the photon-starvation clipping mechanism of \Cref{eq:meth_metal_starvation}. The author also implemented the cascaded U-Net training framework, the dual-filter alternative-input scheme of \Cref{eq:meth_alt_input}, the residual-cascade loss formulation of \Cref{eq:meth_loss_residual}, the baseline pipelines (\gls{bm3d}, single-stage U-Net, and RED-CNN), and the full experimental and reporting workflow. The U-Net backbone follows \cite{Ronneberger2015}, while the residual-prediction principle follows prior work such as DnCNN \cite{Zhang2017} and ResNet \cite{He2016}; in this thesis these established ideas are adapted to the specific cascade comparison studied here. The LoDoPaB-CT dataset and its inherited low-dose sinograms are due to Leuschner et al.~\cite{Leuschner2021}.

\section{Limitations}
\label{sec:conc_limitations}

Four limitations bound the conclusions of this thesis.

\paragraph{Single dataset and simulated corruption.}
All experiments were conducted on \gls{lodopab}, whose ground-truth images are themselves reconstructed from clinical scan data rather than acquired as physical phantom references, and whose low-dose Poisson noise was simulated at $N_0=\num{4096}$. The three physical artifacts considered here were also injected synthetically. Although the artifact models are grounded in the physical literature~\cite{LuMackie2002, Mao2007, Sijbers2004, Prell2009, Gjesteby2016}, the present protocol does not capture effects such as polyenergetic beam hardening, scatter, or non-rigid motion. The ceiling result should therefore be re-tested on projection data that are closer to clinical acquisition.

\paragraph{Independent artifact sampling.}
The three artifact classes are gated independently through Bernoulli sampling in \Cref{eq:meth_compound}. Clinically meaningful correlations between artifact types are therefore not modeled. A conditional sampler informed by scan context or indication would provide a more realistic corruption distribution.

\paragraph{Cascade depth limited to two.}
The experiments in this thesis are limited to two-stage cascades. It therefore remains open whether deeper cascades would provide additional benefit if new information were injected at each stage. The gradient attenuation observed in \Cref{sec:exp_gradient} suggests that naive end-to-end scaling is unlikely to help under the present protocol, but a controlled three-stage study would be needed to test this directly.

\paragraph{Image-domain scope.}
The architectural conclusion reached here applies specifically to cascades that operate in the image domain after filtered backprojection. It should not be generalized automatically to dual-domain systems, learned reconstruction operators, or iterative methods that can exploit sinogram information unavailable to the present models.

\section{Future Work}
\label{sec:conc_future}

The most natural next step is to break the effective image-domain ceiling identified in \Cref{sec:exp_discussion} by supplying later stages with genuinely new information. Two directions emerge directly from the present results.

\paragraph{Sinogram-domain refinement.}
A lightweight network operating on the corrupted sinogram $\tilde{\mathbf{s}}$, followed by reconstruction and fusion with the image-domain Stage~1 output, would give Stage~2 access to information that is not derivable from $\hat{\mathbf{x}}_1$ alone. This is the natural extension suggested by the dual-filter result: the second stage became more useful precisely when it was given an alternative view of the same corrupted measurement.

\paragraph{Projection-domain validation on other datasets.}
Repeating the cascade study on projection-domain datasets with different scanner geometries and anatomy would test whether the image-domain ceiling observed here is specific to LoDoPaB-CT or reflects a more general limitation of post-reconstruction restoration under compound corruption.

\paragraph{Learned reconstruction operators.}
A complementary direction is to relax the fixed FBP operator and allow part of the reconstruction step itself to be learned, for example through learned filtered backprojection or a lightly unrolled iterative scheme. This would allow the model to redistribute capacity across reconstruction and restoration rather than restricting all adaptation to post-reconstruction image processing.

\paragraph{Beyond two stages with information injection.}
A controlled study of three-stage cascades, with additional information injected at every stage, would help distinguish whether the observed ceiling is fundamentally per-stage or per-cascade. On the basis of the present results, residual coupling with explicit information injection would be the most promising formulation to scale beyond two stages.